\documentclass[10pt,twocolumn]{article}

% ---------------------------------------------------------------------------
% Geometry-Conditioned, Repair-Guaranteed Zero-Cost NAS for Unseen Data
% Self-contained: compiles with a vanilla pdflatex install (no external .sty).
% ---------------------------------------------------------------------------
\usepackage[utf8]{inputenc}
\usepackage[T1]{fontenc}
\usepackage{lmodern}
\usepackage[margin=0.85in]{geometry}
\usepackage{amsmath,amssymb}
\usepackage{graphicx}
\usepackage{booktabs}
\usepackage{algorithm}
\usepackage{algpseudocode}
\usepackage{listings}
\usepackage{xcolor}
\usepackage{url}
\usepackage[hidelinks]{hyperref}
\usepackage{caption}
\usepackage{microtype}

% ---- listings style for genotype / repair pseudocode -----------------------
\definecolor{codebg}{rgb}{0.97,0.97,0.97}
\definecolor{codekw}{rgb}{0.10,0.10,0.55}
\definecolor{codecm}{rgb}{0.35,0.50,0.35}
\lstset{
  basicstyle=\ttfamily\scriptsize,
  backgroundcolor=\color{codebg},
  keywordstyle=\color{codekw}\bfseries,
  commentstyle=\color{codecm}\itshape,
  breaklines=true,
  frame=single,
  framesep=3pt,
  columns=fullflexible,
  showstringspaces=false,
  language=Python,
}

\newcommand{\family}[1]{\textsc{#1}}
\newcommand{\code}[1]{\texttt{\small #1}}

\title{\vspace{-1.5em}\textbf{Geometry-Conditioned, Repair-Guaranteed Zero-Cost
Neural Architecture Search for Truly Unseen Data}}

\author{
NAS Unseen-Data 2026 Submission\\
\small Technical Report
}
\date{}

\begin{document}
\maketitle

\begin{abstract}
We present a neural architecture search (NAS) system designed for the
\emph{unseen-data} regime, in which the search algorithm is given only a tensor
of unknown semantics --- its modality, label structure, and even spatial
meaning are hidden --- together with a wall-clock budget. Standard NAS
pipelines assume a fixed, hand-tuned search space for a known domain (e.g.\
$32{\times}32$ natural images) and break down when handed a $1{\times}9{\times}9$
symbolic grid or a $1{\times}6{\times}768$ pseudo-sequence. Our central insight is
that, while the \emph{semantics} of an unseen tensor are unknown, its
\emph{geometry} is fully observable and carries strong structural priors. Our
primary contribution is the \emph{system}: the synthesis that makes zero-cost,
evolutionary NAS run unattended and end-to-end across radically different data
geometries with no per-dataset tuning. Within it we contribute: (i) a
deterministic \emph{family inference} map from input geometry to hard
architectural constraints and soft sampling biases; (ii) a
\emph{spatial-aware repair operator}, applied inside the evolutionary loop, that
projects any sampled or mutated genotype onto the feasible set with correct
ceil/floor down-sampling semantics --- turning mutation into a Markov kernel
that never produces a broken architecture; (iii) a \emph{hardware-aware memory
budget} folded into repair so that feasibility accounts for true training-time
memory, not just forward activations; and (iv) a \emph{three-zone
early-stopping} rule that distinguishes regression from plateau and therefore
terminates cleanly even on tasks driven to their theoretical accuracy ceiling.
We adapt the AZ-NAS zero-cost proxies to a free-form, heterogeneous search space
via a gradient-retaining feature hook with per-proxy graceful degradation, and
drive an \emph{anytime} aging-evolution search that scales to whatever
wall-clock budget each dataset is given. We further detail the robustness
engineering --- a three-level defense against shape errors, full seeding under
seed-sensitive zero-cost proxies, and tolerant metadata parsing --- that lets
the system survive thirteen unknown datasets unattended.
\end{abstract}

% ===========================================================================
\section{Introduction}
% ===========================================================================

Neural architecture search has matured into a reliable tool \emph{when the
target domain is known}. A practitioner searching for a CIFAR-10 classifier
chooses a search space whose primitives, resolutions, and down-sampling
schedule already encode decades of vision priors. The NAS Unseen-Data
challenge inverts this assumption: the system receives a tensor of shape
$N\times C\times H\times W$ with no information about what the data
\emph{means}. Across the challenge's practice datasets a single pipeline must
cope with, among others:

\begin{itemize}\itemsep2pt
\item tiny symbolic grids ($1{\times}9{\times}9$ Sudoku / Game-of-Life),
\item multi-channel board states ($12{\times}8{\times}8$ chess tensors),
\item elongated pseudo-sequences ($1{\times}6{\times}768$ cryptic crosswords,
      $1{\times}28{\times}768$ language tensors),
\item volumetric data ($20{\times}20{\times}20$ voxels),
\item and conventional images ($3{\times}28{\times}28$ to $3{\times}128{\times}128$).
\end{itemize}

A search space tuned for any one of these is pathological for the others: a
fixed four-stage stride-2 schedule annihilates a $9{\times}9$ grid after two
stages; a square $3{\times}3$ kernel wastes capacity on a $1{\times}768$ axis;
an attention head that is cheap at $8{\times}8$ is ruinous at
$128{\times}128$. The naive failure modes are not subtle --- they are shape
errors, out-of-memory crashes, and architectures that simply cannot be built.

\paragraph{Thesis.} The semantics of an unseen tensor are hidden, but its
\emph{geometry is not}. The number of channels, the spatial aspect ratio, the
absolute spatial extent, and the number of classes are all observable before
any training, and each is a strong structural prior. Our system turns geometry
into constraints, constraints into a feasible search space, and a feasible
search space into a robust, time-bounded search.

% ===========================================================================
\section{What Is Original Here}
\label{sec:original}
% ===========================================================================

Because this system deliberately reuses strong existing components, we begin by
drawing a sharp line between prior art and our contribution.

\subsection{Reused prior art (not claimed)}
\begin{itemize}\itemsep2pt
\item \textbf{AZ-NAS zero-cost proxies.} The four proxy signals ---
  \emph{expressivity} (eigenvalue entropy of per-layer feature covariance),
  \emph{progressivity} (minimum inter-layer entropy increment),
  \emph{trainability} (inter-layer Jacobian spectral score), and a
  \emph{complexity} penalty --- are taken from AZ-NAS~\cite{aznas2024}.
\item \textbf{Aging (regularized) evolution.} The population search with
  tournament selection and oldest-first eviction follows
  Real et al.~\cite{regevo2019}.
\item \textbf{Block primitives.} Inverted-residual / MBConv~\cite{mobilenetv2},
  squeeze-excitation~\cite{senet}, stochastic depth /
  DropPath~\cite{stochasticdepth}, separable and bottleneck convolutions are
  all standard.
\end{itemize}

\subsection{Original contributions (claimed)}
\label{sec:contributions}

Our primary contribution is a \emph{system}: the particular synthesis of
components below is, to our knowledge, the first to make zero-cost,
evolutionary NAS run \emph{unattended and end-to-end across radically different
data geometries} without any per-dataset human tuning. No single ingredient is
sufficient --- a geometry prior with no repair still crashes; repair with a
fixed search space still wastes capacity; a zero-cost proxy designed for a
fixed cell does not even run on a $9{\times}9$ grid. The contribution is that
these pieces are made mutually compatible and self-correcting. We decompose it
into the following claims.

\paragraph{(A) System-level: an anytime, geometry-to-architecture pipeline for
unseen data.} A complete method that ingests an uninterpreted tensor and a
wall-clock budget and returns a trained classifier, scaling its own
search/train effort to whatever time it is given and never overrunning. The
end-to-end coupling --- geometry $\rightarrow$ constraints $\rightarrow$
repair-guaranteed feasible space $\rightarrow$ proxy-driven evolution
$\rightarrow$ time-budgeted training --- is the headline result.

\paragraph{(B) Component contributions.}
\begin{enumerate}\itemsep3pt
\item \textbf{Geometry$\rightarrow$family inference} (\S\ref{sec:family}). A
  deterministic map \code{infer\_family(C,H,W,n\_cls)} from raw input geometry
  to one of eight \emph{family profiles}, each fixing hard constraints (max
  down-sampling steps, normalization, anisotropy axis, attention
  admissibility) and soft sampling biases. Conditioning the search space itself
  on observable geometry --- rather than on a known dataset identity --- is the
  move that makes the unseen-data setting tractable.

\item \textbf{Repair as a feasibility projection inside the evolutionary loop}
  (\S\ref{sec:repair}). A deterministic projection \code{repair(g,\dots)} that
  maps \emph{any} genotype onto the feasible set by tracking true spatial size
  with correct ceil/floor down-sampling semantics. Applied to both
  initialization and \emph{every} mutation, repair turns the mutation operator
  into a Markov kernel that always lands in the feasible region --- so the
  search never spends a single evaluation on a broken architecture, and the
  effective search distribution is the feasible set, not the raw genotype
  space.

\item \textbf{Hardware-aware memory feasibility} (\S\ref{sec:memory}). A
  training-memory estimator that counts \emph{expanded intermediate}
  activations (which dominate at large spatial extent) under a training
  overhead factor, with a budget tied to the \emph{actual} GPU. The constraint
  moves out of the training loop (where it is a crash) and into the search
  space (where it is, at worst, a slightly smaller still-strong network).

\item \textbf{Three-zone early stopping} (\S\ref{sec:es}). Separate
  \emph{regression} and \emph{plateau} counters that terminate cleanly both on
  worsening runs and on runs pinned at a theoretical-maximum accuracy, where
  standard no-improvement patience never fires.

\item \textbf{Free-form AZ-NAS adaptation with graceful degradation}
  (\S\ref{sec:proxy-adapt}). A gradient-retaining \code{extract\_layer\_features}
  hook lets AZ-NAS proxies --- written for the fixed NAS-Bench-201 cell --- run
  on a heterogeneous, variable-depth space. Crucially, the score is assembled
  from \emph{whichever} proxies remain finite for a given geometry, so the
  \emph{same} proxy survives both a $9{\times}9$ grid and a
  $128{\times}128$ image without special-casing. This per-proxy degradation is
  itself a generalization mechanism, not just error handling.
\end{enumerate}

\paragraph{(C) Design elements that, while individually small, are part of the
novelty and are easy to miss.}
\begin{itemize}\itemsep2pt
\item \textbf{Vestigial-stage genome.} The genome always carries the maximum
  number of stage genes but decodes only the active prefix. Depth mutations are
  therefore reversible and never destroy genetic material, preserving diversity
  under variable-depth evolution.
\item \textbf{The head is a free gene.} We impose \emph{no} prior on the
  classifier head (seven pooling/MLP variants compete), on the principle that
  the right read-out for an unknown geometry is exactly what search should
  discover.
\item \textbf{A coherent anisotropy thread.} The detected anisotropy axis flows
  from family inference through factored $(1{\times}k)+(k{\times}1)$ blocks to
  axis-aligned pooling in the decoder, giving elongated $1{\times}N$ tensors a
  structurally appropriate operator set.
\item \textbf{Defense-in-depth for shape correctness} (\S\ref{sec:defense}):
  block-level spatial invariants, static repair, and a dynamic dry run --- three
  independent layers, each catching what the previous might miss.
\item \textbf{Determinism as a correctness property} (\S\ref{sec:reproducibility}).
  Because zero-cost proxies are RNG-sensitive (trainability draws Rademacher
  vectors via \code{torch.bernoulli}), seeding every source is what makes the
  \emph{ranking}, and hence the search, reproducible rather than merely tidy.
\end{itemize}

The remainder of the paper develops each contribution, then describes the
end-to-end pipeline.

% ===========================================================================
\section{Geometry-Conditioned Family Inference}
\label{sec:family}
% ===========================================================================

The first stage observes only $(C,H,W,\text{n\_cls})$ and emits a
\code{FamilyProfile}. Let $s_{\min}=\min(H,W)$, $s_{\max}=\max(H,W)$,
$A=HW$, and the anisotropy ratio $\rho=s_{\max}/\max(s_{\min},1)$. The rules
are evaluated in priority order; the first match wins.

\begin{table}[t]
\centering
\caption{The eight geometry-derived families. ``Pools'' is the hard cap on
down-sampling stages; ``GN'' forces GroupNorm; ``Aniso'' marks an anisotropic
down-sampling axis. Block lists are \emph{soft} sampling biases, not hard
constraints.}
\label{tab:families}
\footnotesize
\begin{tabular}{@{}llccc@{}}
\toprule
Family & Trigger (priority order) & Pools & GN & Aniso\\
\midrule
\family{anisotropic}        & $\rho\!\ge\!6,\ s_{\min}\!\le\!8$           & $\le2$ & \checkmark & \checkmark\\
\family{small\_grid}        & $s_{\max}\!\le\!10$                          & $1$    & \checkmark & --\\
\family{possible\_voxel}    & $H\!\approx\!W\!\approx\!C,\ A\!\le\!625$     & $2$    & \checkmark & --\\
\family{channel\_heavy}     & $C\!\ge\!8,\ A\!\le\!1024$                    & $2$    & cond.\ & --\\
\family{spatiotemporal}     & $C\!\ge\!3,\ (H{=}1\!\lor\!W{=}1),s_{\max}\!\ge\!32$ & $3$ & -- & \checkmark\\
\family{visual\_large}      & $s_{\min}\!\ge\!64$                          & $5$    & -- & --\\
\family{visual\_medium}     & $s_{\min}\!\ge\!24$                          & $3$    & -- & --\\
\family{compact\_general}   & otherwise                                    & $2$    & -- & --\\
\bottomrule
\end{tabular}
\end{table}

Two design choices deserve emphasis. First, the head type is left a
\emph{fully free gene} --- we deliberately impose \emph{no} prior on the
classifier head, since the relationship between geometry and the right pooling
strategy is exactly what the search should discover. Second, the
\family{small\_grid} and anisotropic families force GroupNorm: a
$9{\times}9$ feature map yields few spatial samples per channel, so
BatchNorm statistics are unreliable, whereas GroupNorm is batch-independent.
These are the kind of priors a human would apply after \emph{seeing} the data;
we derive them from geometry alone.

% ===========================================================================
\section{The Genotype and Its Decoding}
\label{sec:genotype}
% ===========================================================================

An architecture is a fixed-width \emph{genotype}: a small set of global genes
(stem type, stem width, neck, head, head dropout, normalization, and an active
stage count $n\!\le\!5$) plus five \emph{stage genes}. Each stage gene is an
integer-indexed tuple over discrete option lists: block type (11 choices),
kernel $\in\{1,3,5,7\}$, blocks-per-stage $\in\{1,2,3,4\}$, channel width (9
levels from 16 to 256), expansion, dilation, squeeze-excitation toggle and
ratio, skip mode, down-sampling op $\in\{\text{stride2, maxpool, avgpool,
identity}\}$, and drop-path rate. Vestigial stages beyond $n$ are retained in
the genome but ignored at build time, so mutation can grow or shrink depth
without destroying genetic material.

The decoder (\code{build\_model}) is geometry-agnostic by construction: it
chains a stem, the active stages (each optionally preceded by a separate
down-sampling module), a neck, and a head. The decisive design element is the
feature hook described in \S\ref{sec:proxy-adapt}.

% ===========================================================================
\section{Spatial-Aware Repair}
\label{sec:repair}
% ===========================================================================

Random sampling and mutation over the genotype space routinely produce
\emph{infeasible} architectures: a $5{\times}5$ kernel on a $4{\times}4$ map, a
fifth down-sampling stage on an already-$1{\times}1$ tensor, an attention block
on a degenerate spatial grid. Rather than reject-and-resample (wasteful) or
hope-and-crash (fatal), we \emph{repair}: a deterministic operator that
projects any genotype onto the feasible set.

\subsection{The ceil/floor subtlety}
The repair operator's correctness hinges on tracking the true spatial size.
PyTorch's strided convolution with ``same'' padding rounds \emph{up}
($\lceil h/2\rceil$, e.g.\ $9\!\to\!5$), whereas max/avg pooling rounds
\emph{down} ($\lfloor h/2\rfloor$, e.g.\ $9\!\to\!4$). A repair pass that
assumes a single rule produces a model whose main path and skip connection
disagree on spatial size --- the classic ``expected 5 got 4'' runtime error.
We encode the distinction explicitly:

\begin{lstlisting}
def _actual_spatial(h, ds):
    if ds == 'identity': return h
    if ds == 'stride2':  return (h + 1) // 2   # ceil
    return max(1, h // 2)                       # floor (pool)
\end{lstlisting}

\subsection{Eleven ordered rules}
Repair applies, in order: (R1) clamp stage count; (R2) apply the family's
normalization override; (R3--R4) replace family-forbidden and
geometry-illegal blocks; (R5--R7) a single spatial-aware pass that --- per
stage --- converts strided convolution on an odd dimension to pooling (so
floor semantics hold everywhere), enforces the family's pooling budget, clamps
the kernel below the current spatial size, and clamps dilation so the effective
receptive field fits; (R8) enforce monotonically non-decreasing channel widths;
(R9--R10) coerce the head to a geometry-compatible variant; and (R11) the
hardware-aware memory guard of \S\ref{sec:memory}. The output is guaranteed
constructible, and --- in combination with the dry run of
\S\ref{sec:defense} --- guaranteed to execute a forward pass.

% ===========================================================================
\section{Hardware-Aware Memory Feasibility}
\label{sec:memory}
% ===========================================================================

A subtle but critical failure mode is the architecture that builds, scores
well under zero-cost proxies, and \emph{then} exhausts GPU memory during
training. We observed exactly this: a model on a large-spatial input allocated
$>22$\,GB and crashed, producing zero predictions for that dataset.

\subsection{Estimating the right quantity}
The root cause was a memory estimator that counted only stage-\emph{output}
activations. This ignores the two dominant terms during training:

\begin{enumerate}\itemsep2pt
\item \textbf{Expanded intermediates.} An MBConv block expands channels by up
  to $6\times$ internally; at large spatial extent this intermediate tensor,
  not the block output, dominates memory.
\item \textbf{Training overhead.} Backward-pass gradient buffers and the
  Adam optimizer's two moment tensors multiply the activation footprint.
\end{enumerate}

Our estimator therefore sums, over active stages, both output and
\emph{expanded} intermediate activations, and applies a conservative training
factor:
\begin{equation}
\hat{M} \;=\; \tau \cdot 4 \!\!\sum_{\text{stages}}\!\! B\big(c_{\text{out}}hw +
n_{\text{blk}}\,c_{\text{mid}}\,hw\big)\ \text{bytes},
\end{equation}
with $c_{\text{mid}}=\max(c_{\text{in}},c_{\text{out}})\cdot\text{expansion}$,
batch $B$, and training factor $\tau=3$. On the worst-case model that crashed,
this estimate is $21.9$\,GB --- matching the observed allocation, versus the
$\sim0.2$\,GB the old estimator reported.

\subsection{Budget tied to real hardware}
A fixed budget would be wrong in both directions: too small forces needlessly
tiny networks; too large fails to prevent the crash. We instead set the budget
to $80\%$ of the \emph{physical} GPU memory queried at runtime
(\code{torch.cuda.get\_device\_properties}), with a 16\,GB fallback. Because the
estimator is deliberately conservative (it bills activations at fp32 while
training actually runs in fp16 / AMP), $80\%$ of physical memory is a safe
target. The guard then shrinks an over-budget genotype in capacity-preserving
order --- expansion first, then width, then depth --- for up to three rounds.

\begin{table}[t]
\centering
\caption{Memory guard behavior. Only the genuinely GPU-exceeding model is
altered; normal-geometry models pass untouched, preserving the strong
architectures the search finds. Estimates at $B{=}32$.}
\label{tab:memory}
\footnotesize
\begin{tabular}{@{}lrrl@{}}
\toprule
Input geometry & Before & After & Action\\
\midrule
$3{\times}32{\times}32$ (CIFAR-scale)  & 0.6\,GB  & 0.6\,GB  & untouched\\
$3{\times}96{\times}96$                & 5.0\,GB  & 5.0\,GB  & untouched\\
$3{\times}128{\times}128$              & 9.0\,GB  & 9.0\,GB  & untouched\\
$1{\times}160{\times}160$ (worst case) & 21.9\,GB & 10.8\,GB & shrunk\\
\bottomrule
\end{tabular}
\end{table}

Table~\ref{tab:memory} captures the design goal precisely: the guard is
\emph{inert} for all normal geometries and activates only for the catastrophic
case, where it preserves a strong $192$-channel, four-stage network rather than
collapsing to a degenerate one. A trainable smaller network strictly dominates
an out-of-memory crash.

% ===========================================================================
\section{Adapting AZ-NAS to a Free-Form Space}
\label{sec:proxy-adapt}
% ===========================================================================

AZ-NAS proxies require per-layer features with retained gradients --- the
trainability proxy back-propagates a random cotangent between consecutive
layers and scores the resulting Jacobian's spectrum. The original
implementation relies on the fixed NAS-Bench-201 cell interface. We instead
have the decoded model expose a hook:

\begin{lstlisting}
def extract_layer_features(self, x):
    feats = []
    x = self.stem(x)
    x = x.clone().requires_grad_(True); x.retain_grad()
    feats.append(x)
    for stage, ds in zip(self.stages, self.downs):
        if ds is not None: x = ds(x)
        x = stage(x)
        x = x.clone().requires_grad_(True); x.retain_grad()
        feats.append(x)
    return feats
\end{lstlisting}

This makes the identical proxy math apply to a heterogeneous, variable-depth
network. The combined score is
\begin{equation}
S \;=\; E + P + T \;-\; \lambda_c \log(\#\text{params}+1),
\end{equation}
with $\lambda_c=0.1$, after a Kaiming fan-in re-initialization.

\paragraph{Graceful degradation.} Some geometries make a proxy ill-defined ---
e.g.\ a $9\!\to\!4$ down-sampling makes the trainability Jacobian's spatial
strides non-divisible, so the layer pair is skipped. We compute $S$ from
whichever of $\{E,P,T\}$ remain finite, returning $-\infty$ only if all fail.
This single change is what lets the \emph{same} proxy survive Sudoku
($9{\times}9$) and Myofibre ($128{\times}128$) without special-casing.

% ===========================================================================
\section{Aging-Evolution Search}
\label{sec:search}
% ===========================================================================

We search with regularized (aging) evolution~\cite{regevo2019}: maintain a
population, repeatedly draw a tournament, mutate its fittest member, evaluate
the child's zero-cost score, append it, and evict the \emph{oldest} individual.
Age-based eviction keeps the population churning and resists premature
convergence. Three elements are tailored to our setting:

\begin{itemize}\itemsep2pt
\item \textbf{Repair-in-the-loop.} Every sampled and mutated genotype passes
  through \code{repair} \emph{and} a CPU dry run before scoring, so the
  population never contains an infeasible individual.
\item \textbf{Mutation-scale annealing.} The mutation operator anneals from
  \emph{large} (replace a whole stage or a global gene) early, through
  \emph{medium}, to \emph{small} (perturb one field) late, trading exploration
  for exploitation as the budget is consumed.
\item \textbf{Time as the true terminator.} The nominal round count
  ($2000$) is a ceiling; the search actually runs until a fraction of the
  wall-clock budget is spent, so it scales to whatever time the dataset is
  allotted.
\end{itemize}

\begin{algorithm}[t]
\caption{Repair-guarded aging evolution (abridged)}
\label{alg:evo}
\begin{algorithmic}[1]
\State seed all RNGs \Comment{determinism, \S\ref{sec:reproducibility}}
\State $\mathcal{P}\gets\emptyset$
\While{$|\mathcal{P}|<P$ and time left}
  \State $g\gets$ \Call{Repair}{\Call{Sample}{family}}
  \If{\Call{DryRun}{$g$} ok} $\mathcal{P}\gets\mathcal{P}\cup\{(g,\textsc{Proxy}(g))\}$ \EndIf
\EndWhile
\While{round $<R$ and time left}
  \State $t\gets$ tournament of $k$ from $\mathcal{P}$
  \State $g'\gets$ \Call{Repair}{\Call{Mutate}{$\arg\max_{\text{fit}}t$, scale}}
  \State score $g'$ via dry-run-checked \textsc{Proxy}; append
  \State age all; evict oldest
\EndWhile
\State \Return $\arg\max_{\text{fit}}\mathcal{P}$
\end{algorithmic}
\end{algorithm}

% ===========================================================================
\section{Anytime Budgeting}
\label{sec:anytime}
% ===========================================================================

The competition allots each dataset a wall-clock \code{time\_limit}, and a
pipeline that overruns is penalized while one that finishes early wastes
compute. Our pipeline is therefore \emph{anytime} at two scales. Within a
dataset, search consumes a fixed fraction of the budget and training claims the
remainder, but training additionally projects its own per-epoch cost (a moving
average of recent epoch times) and stops \emph{before} an epoch would breach the
deadline --- so it neither overruns nor leaves an epoch's worth of time on the
table. Across datasets, an optional global pool redistributes time saved on an
easy dataset to the remaining ones, subject to a per-dataset floor. Combined
with the time-terminated evolution of \S\ref{sec:search}, the same code runs
sensibly whether a dataset is given thirty minutes or seven hours, with no
hyper-parameter change.

% ===========================================================================
\section{Three-Zone Early Stopping}
\label{sec:es}
% ===========================================================================

Final training is time-budgeted, but we still want to stop a model that has
converged so the saved time returns to the global pool. Standard
``stop after $p$ epochs without improvement'' fails in two opposite regimes we
actually hit: it stops \emph{too early} when a still-learning model oscillates,
and it \emph{never} stops a model pinned at a theoretical accuracy ceiling
(e.g.\ a deterministic task at $100\%$ validation), where the metric neither
improves nor regresses.

We therefore classify each epoch into three zones relative to the running best
$v^\star$, an improvement threshold $\delta$ (which decays as improvements
accumulate), and a regression threshold $\delta_r$:
\[
\text{epoch} \in
\begin{cases}
\text{$\star$\ improvement} & v > v^\star + \delta\\
\text{$\downarrow$ regression} & v < v^\star - \delta_r\\
\text{$\sim$ plateau} & \text{otherwise.}
\end{cases}
\]
Improvement resets both counters and updates $v^\star$; regression increments a
patience counter; plateau increments a \emph{separate} plateau counter.
Training stops when either counter saturates. The plateau counter is the key
addition: it guarantees termination on ceiling-pinned tasks, while the
regression counter tolerates the benign oscillations that wrongly tripped
naive patience. The best weights (by validation accuracy) are always restored
before the final report.

% ===========================================================================
\section{Robustness Engineering}
\label{sec:defense}
% ===========================================================================

A competition pipeline that crashes on one dataset forfeits that dataset
entirely. We hardened the system along three axes.

\paragraph{Three-level defense against shape errors.}
(1) Block constructors keep main-path and skip-path spatial sizes consistent
(e.g.\ the attention block pools its skip with the same floor semantics as its
body); (2) the repair operator of \S\ref{sec:repair} statically prevents
infeasible geometry; (3) a CPU dry run with a batch of two executes a real
forward pass after \code{build\_model} and before any expensive proxy
evaluation, so any residual shape error is caught and the individual silently
discarded rather than aborting the search.

\paragraph{Per-batch OOM recovery.}
As a last resort, the training loop catches
\code{torch.cuda.OutOfMemoryError} per batch, clears the cache, and skips the
batch; three consecutive OOMs escalate to a clear diagnostic rather than a raw
stack trace.

\paragraph{Tolerant metadata parsing.}
Competition metadata files are external data and not always valid JSON.
We encountered both a UTF-8 BOM and \emph{bare unquoted tokens} as values
(e.g.\ \code{"benchmark": ?}), which abort \code{json.loads} at a byte offset
with no obvious cause. Our loader strips the BOM and coerces a family of
non-standard bare tokens --- \code{?}, \code{NA}, \code{N/A}, \code{NaN},
\code{Inf}/\code{-Inf}/\code{Infinity}, \code{None}, \code{undefined}, and
Python-cased \code{True}/\code{False} --- to valid JSON, using look-ahead so
quoted strings like \code{"NA"} are never touched. The sanitized file is
rewritten so the (unmodifiable) competition runner also reads valid JSON.

% ===========================================================================
\section{Determinism Under Seed-Sensitive Proxies}
\label{sec:reproducibility}
% ===========================================================================

Zero-cost proxies are unusually RNG-sensitive: AZ-NAS trainability draws a
Rademacher cotangent via \code{torch.bernoulli}, and the proxy's value ---
hence the relative ranking of two architectures, hence the entire evolutionary
trajectory --- shifts with the random state. Reproducibility is therefore not
cosmetic; it is a correctness property of the search. A single
\code{set\_seeds} call seeds Python's \code{random}, NumPy, and Torch
(CPU+CUDA), and enables deterministic cuDNN, at the head of both search and
training. We verified that two independent runs with the same seed return the
identical best genotype and fitness.

% ===========================================================================
\section{Experimental Protocol}
\label{sec:protocol}
% ===========================================================================

The system is evaluated on the thirteen NAS Unseen-Data practice datasets. For
each, the pipeline receives the data tensors and a per-dataset wall-clock
\code{time\_limit}; it allocates a fraction to search and the remainder to
time-budgeted training, reserving a buffer for prediction. Data processing
derives per-channel normalization from the training split, applies geometry-aware
augmentation (horizontal flip always; reflective random-crop only when
$H\!\ge\!32$), and selects batch size by pixel count to bound memory. All
randomness is seeded (\S\ref{sec:reproducibility}); the test split is never used
for model selection.

\paragraph{Accuracy results.} Per-dataset test accuracy is reported in
Table~\ref{tab:results}. \emph{These cells are placeholders to be filled from
the official evaluation; we do not report fabricated numbers.}

\begin{table}[t]
\centering
\caption{Per-dataset results (to be completed from official evaluation).
Params and search/train split are emitted automatically by the pipeline's
report writer.}
\label{tab:results}
\footnotesize
\begin{tabular}{@{}lcccc@{}}
\toprule
Dataset & Family & Params & Val & Test\\
\midrule
AddNIST     & \family{visual\_medium} & -- & -- & --\\
Language    & \family{spatiotemporal} & -- & -- & --\\
MultNIST    & \family{visual\_medium} & -- & -- & --\\
CIFARTile   & \family{visual\_large}  & -- & -- & --\\
Gutenberg   & \family{spatiotemporal} & -- & -- & --\\
GeoClassing & \family{visual\_medium} & -- & -- & --\\
Chesseract  & \family{channel\_heavy} & -- & -- & --\\
Sudoku      & \family{small\_grid}    & -- & -- & --\\
Voxel       & \family{possible\_voxel}& -- & -- & --\\
Myofibre    & \family{visual\_large}  & -- & -- & --\\
GameOfLife  & \family{small\_grid}    & -- & -- & --\\
Cryptic     & \family{anisotropic}    & -- & -- & --\\
Windspeed   & \family{spatiotemporal} & -- & -- & --\\
\bottomrule
\end{tabular}
\end{table}

% ===========================================================================
\section{Related Work}
% ===========================================================================

\textbf{Zero-cost NAS.} Training-free proxies estimate architecture quality
from a single minibatch~\cite{naswot,zerocost}. AZ-NAS~\cite{aznas2024}
assembles complementary proxies whose \emph{ranking} correlates strongly with
final accuracy; we adopt its signals but transplant them onto a free-form
search space. \textbf{Evolutionary NAS.} Regularized evolution~\cite{regevo2019}
established age-based eviction; we add repair-in-the-loop and mutation-scale
annealing. \textbf{Search-space design.} Most NAS fixes a domain-specific space
a priori; our family inference instead \emph{derives} the space from observable
geometry, which is the defining requirement of the unseen-data regime.

% ===========================================================================
\section{Conclusion}
% ===========================================================================

We described a NAS \emph{system} whose contribution is as much in the synthesis
as in any single part: the organizing principle is that input \emph{geometry},
though not semantics, is observable and richly informative, and the engineering
is what turns that principle into a pipeline that runs unattended across
thirteen unknown datasets. Geometry yields a family of hard constraints; a
spatial-aware repair operator, applied inside the evolutionary loop, turns those
constraints into a feasible search space where every candidate is buildable ---
making mutation a feasibility-preserving kernel; a hardware-aware memory guard
keeps feasibility honest with respect to the real GPU without collapsing to
degenerate networks; and adapted zero-cost proxies with per-proxy graceful
degradation, driven by repair-guarded aging evolution under an anytime time
budget, select the architecture. Robustness engineering --- a three-level
shape-error defense, deterministic search under seed-sensitive proxies, and
tolerant metadata parsing --- makes the difference between a method that works
on a slide and one that survives thirteen unknown datasets unattended. We view
the central lesson as transferable: when the data is unknown but its shape is
not, condition the search space on geometry and make feasibility a guarantee,
not a hope.

% ===========================================================================
\bibliographystyle{plain}
\begin{thebibliography}{9}
\bibitem{aznas2024}
J.~Lee and B.~Ham.
\newblock {AZ-NAS}: Assembling Zero-Cost Proxies for Network Architecture Search.
\newblock In \emph{CVPR}, 2024.

\bibitem{regevo2019}
E.~Real, A.~Aggarwal, Y.~Huang, and Q.~V. Le.
\newblock Regularized Evolution for Image Classifier Architecture Search.
\newblock In \emph{AAAI}, 2019.

\bibitem{naswot}
J.~Mellor, J.~Turner, A.~Storkey, and E.~J. Crowley.
\newblock Neural Architecture Search without Training.
\newblock In \emph{ICML}, 2021.

\bibitem{zerocost}
M.~S. Abdelfattah, A.~Mehrotra, {\L}.~Dudziak, and N.~D. Lane.
\newblock Zero-Cost Proxies for Lightweight NAS.
\newblock In \emph{ICLR}, 2021.

\bibitem{mobilenetv2}
M.~Sandler, A.~Howard, M.~Zhu, A.~Zhmoginov, and L.-C. Chen.
\newblock {MobileNetV2}: Inverted Residuals and Linear Bottlenecks.
\newblock In \emph{CVPR}, 2018.

\bibitem{senet}
J.~Hu, L.~Shen, and G.~Sun.
\newblock Squeeze-and-Excitation Networks.
\newblock In \emph{CVPR}, 2018.

\bibitem{stochasticdepth}
G.~Huang, Y.~Sun, Z.~Liu, D.~Sedra, and K.~Q. Weinberger.
\newblock Deep Networks with Stochastic Depth.
\newblock In \emph{ECCV}, 2016.

\bibitem{nasbench201}
X.~Dong and Y.~Yang.
\newblock {NAS-Bench-201}: Extending the Scope of Reproducible Neural
Architecture Search.
\newblock In \emph{ICLR}, 2020.
\end{thebibliography}

\end{document}
